\documentclass[11pt,a4paper]{article}
\usepackage[T1]{fontenc}
\usepackage[utf8]{inputenc}
\usepackage{amsmath,amssymb,amsthm}
\usepackage[margin=1in]{geometry}
\usepackage[colorlinks=true,linkcolor=blue,urlcolor=blue]{hyperref}
\theoremstyle{plain}
\newtheorem{theorem}{Theorem}
\newtheorem{lemma}[theorem]{Lemma}
\newtheorem{proposition}[theorem]{Proposition}
\newtheorem{corollary}[theorem]{Corollary}
\theoremstyle{definition}
\newtheorem{remark}[theorem]{Remark}

\title{Recovering the unwritten recurrences of the table OEIS A253985}
\author{Adrian Perez Fontelles\\ \small Independent researcher}
\date{2 September 2026}
\begin{document}
\maketitle

\begin{abstract}
OEIS A253985 is a two-dimensional table $T(n,k)$ whose empirical recurrences are, for several of
its lines, given only as an ORDER --- ``k=2: [order 41]'' and the like --- with the
coefficients in a linked file rather than in the entry. Those conjectures can still be settled,
and the recurrences recovered. A line of the table is a fixed-width or fixed-height array
count, hence a walk count on finitely many vertices, hence satisfies some monic recurrence of
order at most that number; Berlekamp--Massey on twice as many exact terms returns the MINIMAL
one. Where its order equals the order the entry states --- and where the same holds after the
first few terms are dropped --- any recurrence of that order the line satisfies has a
characteristic polynomial that is a multiple of the minimal one and of equal degree, hence is
that polynomial. This note settles column 2, column 3.
\end{abstract}

\noindent\small 2020 Mathematics Subject Classification. 05A15, 11B37, 15A18.\normalsize

\section{The table and the unwritten conjectures}

OEIS A253985 is ``T(n,k)=Number of (n+2)X(k+2) 0..1 arrays with every 3X3 subblock diagonal maximum plus antidiagonal maximum nondecreasing horizontally and vertically.''. It is read by antidiagonals and begins
\[
512, 3372, 3372, 21304, 36434, 21304, 136868, 417518,\ \dots
\]
The entry carries, among its empirical claims:
\begin{quote}\small
k=2: [order 41]\\
k=3: [order 68]
\end{quote}
As of the ``Last modified'' line on the live entry (Jul 23 2025 14:23:35, revision 6) these are still
recorded as empirical, and the recurrences themselves are not in the entry text.

\section{Why a line of the table is a walk count}

Fix the index that the line fixes. The entry defines $T(n,k)$ as a number of arrays of a shape
determined by $n$ and $k$, subject to a condition local to a bounded window of consecutive
lines. Substituting the fixed index into the entry's own wording turns the definition into an
ordinary one-parameter array count, and for such a count the admissible windows are the
vertices of a finite digraph and an array is a walk. So the line is a walk count on $S$
vertices, and by the Cayley--Hamilton theorem it satisfies a monic integer recurrence of order
at most $S$.

\section{Recovering the recurrences}

\begin{lemma}\label{lem:bm}
A sequence known to satisfy some linear recurrence of order at most $S$ is determined, with its
minimal recurrence, by its first $2S$ terms; Berlekamp--Massey applied to them returns that
minimal recurrence exactly.
\end{lemma}

\subsection*{Column $k=2$}
Substituting the fixed index into the entry's wording gives an ordinary one-parameter array
count, a walk on $S=436$ vertices. Berlekamp--Massey on $2S=872$ exact terms
returns a minimal recurrence of order $41$ --- the order the entry states --- with
integer coefficients
\[
(c_1,\dots,c_{41})=(13,\ -7,\ -43,\ -860,\ -920,\ 15830,\ -9400,\ -7539,\ -98847,\ 40329,\ 147709,\ 83340,\ 421956,\ -537724,\ 204762,\ -1909934,\ -520990,\ 1344182,\ -2472492,\ 10604787,\ 1923137,\ -6305887,\ -1207655,\ -18037820,\ 3466634,\ 11862724,\ 5340110,\ 12134681,\ -11925689,\ -8047471,\ 1011577,\ -3455262,\ 6530290,\ 1826872,\ -3037352,\ 578112,\ 205440,\ -198848,\ 22784,\ 10496,\ -1024).
\]
so that $a(m)=\sum_{i=1}^{41}c_i\,a(m-i)$ for every $m$. Removing the first
$41+4$ terms and repeating gives order $41$ again, which is what the
identification below needs.

\subsection*{Column $k=3$}
Substituting the fixed index into the entry's wording gives an ordinary one-parameter array
count, a walk on $S=1748$ vertices. Berlekamp--Massey on $2S=3496$ exact terms
returns a minimal recurrence of order $68$ --- the order the entry states --- with
integer coefficients
\[
(c_1,\dots,c_{68})=(28,\ -145,\ -290,\ 3956,\ -34092,\ 509816,\ -9183232,\ 47960426,\ 37779032,\ -678709222,\ 2723203204,\ -16659134876,\ 53539215576,\ -113330978436,\ 67284010856,\ 1656306886479,\ -8030982358532,\ 32699882935471,\ -80824418372722,\ 56284941969440,\ -46956907765996,\ -1251829135370140,\ 7345279239088264,\ -15642954340140224,\ 34991707665592528,\ -30230892226810528,\ -167431433226489088,\ 525110680374564432,\ -858915930555997632,\ 942642794728024464,\ -581448516624828640,\ -1137297384843599680,\ 9956885340665483712,\ -22016001838561733760,\ 21088569848554387200,\ -6175948085185980672,\ -35119055493109299200,\ 87287433883317381120,\ -71152746818372299776,\ -1193184180774223872,\ 79838548171995246592,\ -145197454741129986048,\ 98219181166194999296,\ 46278324649743876096,\ -234303966570154885120,\ 414271824050006786048,\ -338898926865739415552,\ 28683437128719269888,\ 472882511558017548288,\ -1029184145759225774080,\ 928564152237266305024,\ -275833828457276506112,\ -473425081035810406400,\ 1186363267494275186688,\ -1092181383697229938688,\ 355082531685154160640,\ 216498447847941734400,\ -597837456528830889984,\ 552656415911005126656,\ -175101205924628398080,\ -38203910775694688256,\ 110095182308164239360,\ -101657551517087956992,\ 30118816466551701504,\ 2699600113352835072,\ -4693841527254810624,\ 4103905160441364480,\ -1231171548132409344).
\]
so that $a(m)=\sum_{i=1}^{68}c_i\,a(m-i)$ for every $m$. Removing the first
$68+4$ terms and repeating gives order $68$ again, which is what the
identification below needs.

\begin{theorem}
For each line above, the displayed recurrence holds for every index, and it is the recurrence
the entry records.
\end{theorem}

\begin{proof}
It holds by Lemma~\ref{lem:bm}. For the identification, the entry asserts that the line
satisfies SOME recurrence of the stated order, possibly only from some index on. Let $q$ be its
characteristic polynomial and $q_{\min}$ the minimal polynomial of the tail. Then
$q_{\min}\mid q$, and the second Berlekamp--Massey run --- on the line with its first few
terms removed --- shows $\deg q_{\min}$ equals the stated order, which is $\deg q$. Both are
monic, so $q=q_{\min}$.
\end{proof}

\section{Verification}

Three checks.

First, each line's model was compared against the entry's own published values for that line,
taken from the antidiagonal DATA read in either orientation and from the ``Table starts''
block, and agrees at every available term. This is the only point at which reading the entry's
English enters, and a misreading gives different counts.

Second, each recovered recurrence was evaluated directly on the model's values, with no
Berlekamp--Massey involved, and holds at every index tested.

Third, each order was computed twice by independent routes --- modulo a $61$-bit prime, where
every intermediate is a machine word, and again in exact rational arithmetic --- and once more
on the line with its first few terms dropped, which is what the identification requires. All
agree.

\begin{thebibliography}{9}
\bibitem{oeis} The OEIS Foundation, \emph{The On-Line Encyclopedia of Integer
Sequences}, \url{https://oeis.org}, 2026. Sequence A253985.
\bibitem{massey} J.~L.~Massey, Shift-register synthesis and BCH decoding, \emph{IEEE
Trans. Inform. Theory} \textbf{15} (1969), 122--127.
\bibitem{stanley} R.~P.~Stanley, \emph{Enumerative Combinatorics}, Volume 1, 2nd ed.,
Cambridge University Press, 2011. (Transfer-matrix method, Section 4.7.)
\bibitem{horn} R.~A.~Horn and C.~R.~Johnson, \emph{Matrix Analysis}, 2nd ed., Cambridge
University Press, 2013.
\end{thebibliography}

\end{document}
